\chapter{Complete Grammar}
\label{ch:grammar}

This chapter provides the complete formal grammar for \haira{} in Extended Backus-Naur Form (EBNF), including all core language constructs and agentic extensions (providers, tools, agents, and workflows).

\section{Notation}

\begin{table}[h]
\centering
\begin{tabular}{ll}
\toprule
\textbf{Notation} & \textbf{Meaning} \\
\midrule
\code{=} & Definition \\
\code{|}  & Alternative \\
\code{[ ... ]} & Optional (0 or 1) \\
\code{\{ ... \}} & Repetition (0 or more) \\
\code{( ... )} & Grouping \\
\code{"..."} & Terminal string \\
\code{;} & End of production \\
\bottomrule
\end{tabular}
\caption{EBNF notation}
\end{table}

\section{Lexical Grammar}

\subsection{Characters}

\begin{lstlisting}[language=EBNF]
letter        = "a".."z" | "A".."Z" | "_" ;
digit         = "0".."9" ;
hex_digit     = digit | "a".."f" | "A".."F" ;
octal_digit   = "0".."7" ;
binary_digit  = "0" | "1" ;
any_char      = (* any Unicode character *) ;
\end{lstlisting}

\subsection{Whitespace and Comments}

\begin{lstlisting}[language=EBNF]
whitespace    = " " | "\t" | "\r" ;
newline       = "\n" ;
line_comment  = "//" { any_char } newline ;
block_comment = "/*" { any_char | block_comment } "*/" ;
doc_comment   = "///" { any_char } newline ;
\end{lstlisting}

\begin{note}
Block comments nest: \code{/* outer /* inner */ still outer */} is valid. Doc comments (\code{///}) attach to the immediately following declaration and are preserved in compiled metadata.
\end{note}

\subsection{Identifiers}

\begin{lstlisting}[language=EBNF]
identifier    = letter { letter | digit } ;
\end{lstlisting}

\subsection{Keywords}

\begin{lstlisting}[language=EBNF]
keyword       = "agent" | "and" | "as" | "break" | "catch"
              | "chan" | "const" | "continue" | "defer"
              | "else" | "enum" | "false" | "fn" | "for"
              | "from" | "if" | "impl" | "import" | "in"
              | "match" | "nil" | "not" | "or" | "provider"
              | "return" | "select" | "spawn" | "step" | "struct"
              | "tool" | "true" | "try" | "type" | "unsafe"
              | "where" | "while" | "workflow" ;
\end{lstlisting}

\begin{note}
The keywords \keyword{agent}, \keyword{provider}, \keyword{tool}, and \keyword{workflow} are reserved for agentic declarations. The keyword \keyword{unsafe} is reserved for future use. All keywords are case-sensitive.
\end{note}

\subsection{Literals}

\begin{lstlisting}[language=EBNF]
integer_lit   = decimal_lit | hex_lit | octal_lit | binary_lit ;
decimal_lit   = digit { digit | "_" } ;
hex_lit       = "0" ("x"|"X") hex_digit { hex_digit | "_" } ;
octal_lit     = "0" ("o"|"O") octal_digit { octal_digit | "_" } ;
binary_lit    = "0" ("b"|"B") binary_digit { binary_digit | "_" } ;

float_lit     = decimal_lit "." decimal_lit [ exponent ]
              | decimal_lit exponent ;
exponent      = ("e"|"E") ["+"|"-"] decimal_lit ;

string_lit    = '"' { string_char } '"'
              | "r" '"' { raw_char } '"'
              | triple_string ;
triple_string = '"""' { any_char } '"""' ;
string_char   = any_char - ('"' | "\" | newline)
              | escape_seq ;
escape_seq    = "\" ( "n" | "r" | "t" | "\" | '"' | "0"
              | "x" hex_digit hex_digit
              | "u" "{" hex_digit { hex_digit } "}" ) ;
raw_char      = any_char - '"' ;

bool_lit      = "true" | "false" ;
nil_lit       = "nil" ;
\end{lstlisting}

\begin{note}
Triple-quoted strings (\code{"""..."""}) preserve newlines and leading indentation. They are used for tool descriptions, agent system prompts, and multi-line string literals. The common leading whitespace is stripped at compile time.
\end{note}

\subsection{Operators}

\begin{lstlisting}[language=EBNF]
operator      = arith_op | comp_op | logic_op | bitwise_op
              | shift_op | assign_op | other_op ;

arith_op      = "+" | "-" | "*" | "/" | "%" ;
comp_op       = "==" | "!=" | "<" | ">" | "<=" | ">=" ;
logic_op      = "&&" | "||" | "!" ;
bitwise_op    = "&" | "|" | "^" | "~" ;
shift_op      = "<<" | ">>" ;
assign_op     = "=" | "+=" | "-=" | "*=" | "/="
              | "&=" | "|=" | "^=" | "<<=" | ">>=" | ">>>=" ;
other_op      = "|>" | ".." | "..=" | "<-" | "->" | "=>" | "?" | "@" ;
\end{lstlisting}

\begin{note}
The \code{|>} operator is the pipe operator for function composition. The \code{@} symbol is used for decorator syntax on workflow declarations. The \code{<-} operator is used for channel send and receive operations.
\end{note}

\subsection{Delimiters}

\begin{lstlisting}[language=EBNF]
delimiter     = "(" | ")" | "[" | "]" | "{" | "}"
              | "," | ":" | "." | ";" | "#" ;
\end{lstlisting}

\section{Syntactic Grammar}

\subsection{Program Structure}

\begin{lstlisting}[language=EBNF]
program       = { import_stmt } { top_level } ;

top_level     = fn_decl
              | struct_decl
              | enum_decl
              | type_alias
              | const_decl
              | impl_block
              | constraint_decl
              | provider_decl
              | tool_decl
              | agent_decl
              | workflow_decl ;
\end{lstlisting}

\subsection{Imports}

\begin{lstlisting}[language=EBNF]
import_stmt   = "import" import_spec ;

import_spec   = string_lit
              | identifier "from" string_lit
              | "{" ident_list "}" "from" string_lit
              | "*" "from" string_lit ;

ident_list    = identifier { "," identifier } ;
\end{lstlisting}

\subsection{Core Declarations}

\begin{lstlisting}[language=EBNF]
fn_decl       = [ attributes ] "fn" [ receiver ] identifier
                [ generic_params ] "(" [ params ] ")"
                [ "->" type ] [ where_clause ] block ;

receiver      = "(" identifier ":" [ "*" ] type ")" ;

generic_params = "<" generic_param { "," generic_param } ">" ;

generic_param  = identifier [ ":" constraint_list ] ;

constraint_list = identifier { "+" identifier } ;

where_clause  = "where" where_bound { "," where_bound } ;

where_bound   = identifier ":" constraint_list ;

params        = param { "," param } ;

param         = identifier ":" [ "..." ] type [ "=" expression ] ;

struct_decl   = [ attributes ] "struct" identifier
                [ generic_params ] "{" { field_decl } "}" ;

field_decl    = [ attributes ] identifier ":" type ;

enum_decl     = [ attributes ] "enum" identifier
                [ generic_params ] "{" enum_variants "}" ;

enum_variants = enum_variant { "," enum_variant } [ "," ] ;

enum_variant  = identifier [ "(" type_list ")" ] ;

type_alias    = "type" identifier [ generic_params ] "=" type ;

const_decl    = "const" identifier [ ":" type ] "=" expression ;

impl_block    = "impl" [ generic_params ] identifier
                "for" type [ where_clause ]
                "{" { fn_decl } "}" ;

constraint_decl = "constraint" identifier [ generic_params ]
                  "{" { fn_signature } "}" ;

fn_signature  = "fn" identifier "(" [ type_params ] ")" [ "->" type ] ;

type_params   = type { "," type } ;
\end{lstlisting}

\subsection{Attributes}

Attributes use the \code{\#[...]} syntax and can be applied to functions, structs, enums, and fields. They are distinct from the \code{@} decorator syntax used for workflow triggers.

\begin{lstlisting}[language=EBNF]
attributes    = { attribute } ;

attribute     = "#[" identifier [ "(" attr_args ")" ] "]" ;

attr_args     = attr_arg { "," attr_arg } ;

attr_arg      = identifier [ "=" literal ] ;
\end{lstlisting}

\subsection{Decorators}

Decorators use the \code{@} syntax and are used exclusively for workflow trigger declarations.

\begin{lstlisting}[language=EBNF]
decorator     = "@" identifier [ "(" [ arguments ] ")" ] ;
\end{lstlisting}

\subsection{Provider Declarations}

Providers configure LLM backends. They are declared at the top level and referenced by name in agent declarations.

\begin{lstlisting}[language=EBNF]
provider_decl  = "provider" identifier "{" { provider_field } "}" ;

provider_field = identifier ":" expression ;
\end{lstlisting}

\subsection{Tool Declarations}

Tools are typed functions with a mandatory triple-quoted description. The compiler auto-generates JSON schemas from the type signature.

\begin{lstlisting}[language=EBNF]
tool_decl     = "tool" identifier "(" [ params ] ")" [ "->" type ]
                "{" triple_string { statement } "}" ;

triple_string = '"""' { any_char } '"""' ;
\end{lstlisting}

\begin{note}
The \code{triple\_string} must be the first element of the tool body. The compiler emits an error if a tool is missing its description.
\end{note}

\subsection{Agent Declarations}

Agents are LLM-powered entities with a model, system prompt, optional tools, and optional memory.

\begin{lstlisting}[language=EBNF]
agent_decl    = "agent" identifier "{" { agent_field } "}" ;

agent_field   = identifier ":" agent_value ;

agent_value   = expression
              | string_lit
              | triple_string
              | "[" ident_list "]"
              | memory_value ;

memory_value  = "conversation" "(" [ arguments ] ")"
              | "summary" "(" [ arguments ] ")"
              | "none" ;
\end{lstlisting}

\begin{note}
Agent fields include \code{model} (required), \code{system} (required), and optional fields such as \code{tools}, \code{temperature}, \code{max\_tokens}, \code{max\_steps}, and \code{memory}. The \code{model} field references a provider by name.
\end{note}

\subsection{Workflow Declarations}

Workflows are composable functions with optional trigger decorators. A workflow without a decorator can only be called as a sub-workflow.

\begin{lstlisting}[language=EBNF]
workflow_decl = { decorator } "workflow" identifier
                "(" [ params ] ")" [ "->" type ] block ;
\end{lstlisting}

\subsection{Types}

\begin{lstlisting}[language=EBNF]
type          = simple_type
              | array_type
              | map_type
              | tuple_type
              | option_type
              | fn_type
              | generic_type
              | stream_type
              | pointer_type
              | union_type
              | channel_type ;

simple_type   = "int" | "i8" | "i16" | "i32" | "i64"
              | "u8" | "u16" | "u32" | "u64"
              | "float" | "f32" | "f64"
              | "bool" | "string"
              | identifier ;

array_type    = "[" "]" type ;

map_type      = "[" type ":" type "]" ;

tuple_type    = "(" type { "," type } ")" ;

option_type   = type "?" ;

fn_type       = "fn" "(" [ type_list ] ")" [ "->" type ] ;

type_list     = type { "," type } ;

generic_type  = identifier "<" type_list ">" ;

stream_type   = "stream" "<" type ">" ;

pointer_type  = "*" type ;

union_type    = type { "|" type } ;

channel_type  = "chan" "<" type ">" ;
\end{lstlisting}

\begin{note}
The \type{stream<T>} type represents a lazy, asynchronous, single-use sequence of values. It is distinct from \type{chan<T>}, which is a bidirectional communication channel between concurrent tasks. Streams are typically produced by agent \code{.stream()} calls and consumed with \keyword{for} loops.
\end{note}

\subsection{Statements}

\begin{lstlisting}[language=EBNF]
statement     = var_decl
              | assignment
              | expr_stmt
              | if_stmt
              | for_stmt
              | while_stmt
              | match_stmt
              | return_stmt
              | break_stmt
              | continue_stmt
              | defer_stmt
              | spawn_stmt
              | step_stmt
              | send_stmt
              | select_stmt
              | block ;

var_decl      = ident_list [ ":" type ] "=" expr_list ;

ident_list    = identifier { "," identifier } ;

expr_list     = expression { "," expression } ;

assignment    = lvalue_list assign_op expr_list ;

lvalue_list   = lvalue { "," lvalue } ;

lvalue        = identifier
              | lvalue "." identifier
              | lvalue "[" expression "]" ;

expr_stmt     = expression ;

if_stmt       = "if" [ simple_stmt ";" ] expression block
                [ "else" ( if_stmt | block ) ] ;

simple_stmt   = var_decl | assignment | expr_stmt ;

for_stmt      = "for" [ identifier "," ] identifier "in" expression block ;

while_stmt    = "while" expression block ;

match_stmt    = "match" expression "{" { match_arm } "}" ;

match_arm     = pattern [ "if" expression ] "=>" ( expression | block ) ;

return_stmt   = "return" [ expr_list ] ;

break_stmt    = "break" [ identifier ] ;

continue_stmt = "continue" [ identifier ] ;

defer_stmt    = "defer" statement ;

spawn_stmt    = "spawn" ( block | fn_call ) ;

step_stmt     = "step" string_lit block ;

send_stmt     = expression "<-" expression ;

select_stmt   = "select" "{" { select_case } [ default_case ] "}" ;

select_case   = pattern "=" "<-" expression "=>" ( expression | block )
              | expression "<-" expression "=>" ( expression | block ) ;

default_case  = "default" "=>" ( expression | block ) ;

block         = "{" { statement } "}" ;

fn_call       = identifier "(" [ arguments ] ")"
              | identifier "." identifier "(" [ arguments ] ")" ;
\end{lstlisting}

\begin{note}
Inside a workflow, a \keyword{spawn} block runs all its statements concurrently and returns a list of results when all complete. This is equivalent to \code{Promise.all()} in JavaScript. Outside a workflow, \keyword{spawn} launches a concurrent goroutine-style task. A \keyword{step} block is a named region with automatic telemetry; it is only valid inside \keyword{workflow} bodies. Variables assigned inside a step are visible in subsequent steps and the enclosing workflow scope.
\end{note}

\subsection{Expressions}

\begin{lstlisting}[language=EBNF]
expression    = assign_expr ;

assign_expr   = or_expr [ ("=" | "+=" | "-=" | "*=" | "/="
              | "&=" | "|=" | "^=" | "<<=" | ">>=" | ">>>=")
                assign_expr ] ;

or_expr       = and_expr { ("or" | "||") and_expr } ;

and_expr      = equality { ("and" | "&&") equality } ;

equality      = comparison { ("==" | "!=") comparison } ;

comparison    = pipe { ("<" | ">" | "<=" | ">=") pipe } ;

pipe          = range { "|>" range } ;

range         = bitwise_or [ (".." | "..=") bitwise_or ] ;

bitwise_or    = bitwise_xor { "|" bitwise_xor } ;

bitwise_xor   = bitwise_and { "^" bitwise_and } ;

bitwise_and   = shift { "&" shift } ;

shift         = additive { ("<<" | ">>") additive } ;

additive      = multiplicative { ("+" | "-") multiplicative } ;

multiplicative = unary { ("*" | "/" | "%") unary } ;

unary         = ("!" | "-" | "~" | "not") unary
              | postfix ;

postfix       = primary { postfix_op } ;

postfix_op    = call
              | index
              | field
              | optional_chain ;

call          = "(" [ arguments ] ")" ;

arguments     = argument { "," argument } ;

argument      = [ identifier ":" ] expression ;

index         = "[" expression "]" ;

field         = "." identifier ;

optional_chain = "?" "." identifier ;

primary       = identifier
              | literal
              | "(" expression ")"
              | array_lit
              | map_lit
              | struct_lit
              | closure
              | if_expr
              | match_expr
              | receive_expr
              | block ;

literal       = integer_lit | float_lit | string_lit
              | bool_lit | nil_lit ;

array_lit     = "[" [ expr_list ] "]" ;

map_lit       = "[" ":" "]"
              | "[" map_entry { "," map_entry } [ "," ] "]" ;

map_entry     = expression ":" expression ;

struct_lit    = identifier "{" [ field_inits ] "}" ;

field_inits   = field_init { "," field_init } [ "," ] ;

field_init    = identifier [ ":" expression ] ;

closure       = "fn" "(" [ closure_params ] ")" [ "->" type ]
                ( block | "{" expression "}" ) ;

closure_params = closure_param { "," closure_param } ;

closure_param  = identifier [ ":" type ] ;

if_expr       = "if" expression block "else" ( if_expr | block ) ;

match_expr    = "match" expression "{" { match_arm } "}" ;

receive_expr  = "<-" expression ;
\end{lstlisting}

\subsection{Patterns}

\begin{lstlisting}[language=EBNF]
pattern       = literal_pat
              | ident_pat
              | wildcard_pat
              | tuple_pat
              | struct_pat
              | enum_pat
              | array_pat
              | range_pat
              | or_pat ;

literal_pat   = integer_lit | float_lit | string_lit | bool_lit ;

ident_pat     = identifier ;

wildcard_pat  = "_" ;

tuple_pat     = "(" pattern { "," pattern } ")" ;

struct_pat    = identifier "{" [ field_pats ] "}" ;

field_pats    = field_pat { "," field_pat } ;

field_pat     = identifier [ ":" pattern ] ;

enum_pat      = identifier "." identifier [ "(" [ pattern_list ] ")" ] ;

pattern_list  = pattern { "," pattern } ;

array_pat     = "[" [ pattern_list ] [ "," "..." [ identifier ] ] "]" ;

range_pat     = literal ".." [ "=" ] literal ;

or_pat        = pattern { "|" pattern } ;
\end{lstlisting}

\section{Grammar Summary}

\subsection{Entry Points}

\begin{itemize}
    \item \code{program} --- Complete source file
    \item \code{expression} --- Single expression (REPL)
    \item \code{statement} --- Single statement
    \item \code{type} --- Type annotation
\end{itemize}

\subsection{Precedence (High to Low)}

\begin{enumerate}
    \item Postfix: \code{()} \code{[]} \code{.} \code{?.}
    \item Unary: \code{!} \code{-} \code{\~{}} \code{not}
    \item Multiplicative: \code{*} \code{/} \code{\%}
    \item Additive: \code{+} \code{-}
    \item Shift: \code{<<} \code{>>} \code{>>>}
    \item Bitwise AND: \code{\&}
    \item Bitwise XOR: \code{\^{}}
    \item Bitwise OR: \code{|}
    \item Range: \code{..} \code{..=}
    \item Pipe: \code{|>}
    \item Comparison: \code{<} \code{>} \code{<=} \code{>=}
    \item Equality: \code{==} \code{!=}
    \item Logical AND: \code{and} \code{\&\&}
    \item Logical OR: \code{or} \code{||}
    \item Assignment: \code{=} \code{+=} \code{-=} \code{*=} \code{/=} \code{\&=} \code{|=} \code{\^{}=} \code{<<=} \code{>>=} \code{>>>=}
\end{enumerate}

\subsection{Associativity}

\begin{itemize}
    \item Left-associative: All binary operators except assignment
    \item Right-associative: Assignment operators, unary operators
    \item Non-associative: Comparison operators, range operators
\end{itemize}
